\section*{Capítulo \ref{ch:b2:equilibrium-shifts} --- Constantes de equilibrio y desplazamientos}

\begin{solution}{exo:b2:equilibrium-shifts:1}
$K^\circ = \exp(32\,800/(8.314 \times 298.15)) = \num{5.6e5}$ y, para
$+\qty{32.8}{kJ/mol}$, su inversa, \num{1.8e-6}. Un factor 100 exige
$\Delta(\Delta_r G^\circ) = -RT\ln 100 = \qty{-11.4}{kJ/mol}$.
\end{solution}

\begin{solution}{exo:b2:equilibrium-shifts:2}
$\ln K^\circ(400) = \ln(5.4 \times 10^5) + (91\,880/8.314)(1/400 -
1/298.15) = 13.20 - 9.44 = 3.76$, $K^\circ \approx 43$. Las tablas dan 35.6:
en \qty{100}{K} la aproximación de $\Delta_r H^\circ$ constante ya se
desvía un 20\,\% (la reacción se vuelve más \omterm{def:b2:reaction-enthalpy:exothermic}{exotérmica} a medida que $T$ aumenta).
\end{solution}

\begin{solution}{exo:b2:equilibrium-shifts:3}
(a) $n = 1$, $r = 0$, $\varphi = 2$: $v = 1$ (la presión de vapor es función
de $T$). (b) $v = 3 - 1 + 2 - 3 = 1$. (c) $v = 3 - 1 + 2 - 1 = 3$. (d) Cuatro
especies, una reacción, dos fases (el carbono y el gas): $v = 4 - 1 + 2 - 2 =
3$; si el gas procede solo del carbono y del vapor de agua, $n(\ce{CO}) = n(\ce{H2})$
en el gas, $s = 1$ y $v = 2$.
\end{solution}

\begin{solution}{exo:b2:equilibrium-shifts:4}
Calentar: sentido \omterm{def:b2:reaction-enthalpy:exothermic}{endotérmico}, hacia \ce{NO2} (sentido directo). Comprimir:
hacia menos moléculas de gas, \ce{N2O4} (sentido inverso). Argón a volumen constante:
ningún cambio, pues las presiones parciales de los gases que reaccionan no varían.
\end{solution}

\begin{solution}{exo:b2:equilibrium-shifts:5}
A \qty{298}{K}: $\Delta_r G^\circ = 180.6 - 298.15 \times 0.0248 =
\qty{173.2}{kJ/mol}$, $K^\circ = \num{4.5e-31}$. A \qty{2000}{K}: $\Delta_r
G^\circ = 180.6 - 2000 \times 0.0248 = \qty{131.0}{kJ/mol}$, $K^\circ =
\num{3.8e-4}$. Con $\Delta\nu_{\text{gas}} = 0$, $x(\ce{NO})^2 = K^\circ x(\ce{N2})
x(\ce{O2})$: $x(\ce{NO}) = \sqrt{3.8 \times 10^{-4} \times 0.78 \times 0.21} =
\num{7.9e-3}$, casi un uno por ciento. Los motores queman a esas temperaturas; al
enfriarse los gases, $K^\circ$ cae enormemente, pero la descomposición del \ce{NO} es
muy lenta: el equilibrio en caliente queda congelado en los gases de escape (hasta que un catalizador
lo elimina).
\end{solution}

\begin{solution}{exo:b2:equilibrium-shifts:6}
Con $\xi$ mol de éster y una cantidad total constante, $Q =
\xi^2/[(1 - \xi)(n_B - \xi)]$. Para $n_B = 1.0$: $\xi^2/(1 - \xi)^2 = 4$, $\xi
= \qty{0.67}{mol}$. Para $n_B = 3.0$: $3\xi^2 - 16\xi + 12 = 0$, $\xi =
\qty{0.90}{mol}$. Eliminar el agua mantiene $Q$ por debajo de $K^\circ$ sea cual sea $\xi$: la
reacción sigue avanzando en sentido directo hasta que se consume el ácido.
\end{solution}

\begin{solution}{exo:b2:equilibrium-shifts:7}
Sin argón: cantidades $1 - \alpha$, $2\alpha$, total $1 + \alpha$;
$Q = 4\alpha^2/(1 - \alpha^2) = 0.148$, $\alpha = \sqrt{0.148/4.148} = 0.189$.
Con \qty{1}{mol} de argón a \qty{1}{bar}: total $2 + \alpha$, $Q =
4\alpha^2/[(1 - \alpha)(2 + \alpha)] = 0.148$, luego $4.148\alpha^2 + 0.148\alpha
- 0.296 = 0$, $\alpha = 0.250$: diluir a presión total constante rebaja las
presiones parciales, lo que favorece el lado con más moléculas de gas. A
volumen constante, las presiones parciales de los gases que reaccionan no cambian:
$\alpha$ sigue siendo 0.189.
\end{solution}

\begin{solution}{exo:b2:equilibrium-shifts:8}
Solo el sólido: $n = 3$, $r = 1$, $s = 1$ ($p(\ce{NH3}) = p(\ce{HCl})$),
$\varphi = 2$: $v = 3 - 1 - 1 + 2 - 2 = 1$; se elige $T$ y la presión queda
fijada (una ``presión de disociación''). Con amoníaco añadido, $s = 0$ y $v = 2$:
pueden elegirse $T$ y una presión; el producto $p(\ce{NH3})\,p(\ce{HCl})$ queda
fijado por $T$.
\end{solution}

\begin{solution}{exo:b2:equilibrium-shifts:9}
$\Delta_r H^\circ = -R\ln(K_2/K_1)/(1/T_2 - 1/T_1) = -8.314 \times
\ln(0.0173)/(1/600 - 1/500) = \qty{-101}{kJ/mol}$, más negativa que a
\qty{298}{K} (\qty{-91.9}{kJ/mol}), como predecía la \omterm{thm:b2:reaction-enthalpy:kirchhoff}{ley de Kirchhoff}
($\Delta_r C_p^\circ < 0$).
\end{solution}

\begin{solution}{exo:b2:equilibrium-shifts:10}
$Q = \dfrac{n_{\ce{NH3}}^2\,n_{\text{tot}}^2}{n_{\ce{N2}}n_{\ce{H2}}^3}
\left(\dfrac{p^\circ}{p}\right)^2$. A $T$, $p$ constantes:
$\partial\ln Q/\partial n_{\ce{N2}} = -1/n_{\ce{N2}} + 2/n_{\text{tot}}$, positiva
cuando $n_{\ce{N2}} > n_{\text{tot}}/2$, es decir, $x(\ce{N2}) > 1/2$. Entonces $Q$
supera a $K^\circ$ y el equilibrio se desplaza en sentido inverso. A volumen constante,
$\partial\ln Q/\partial n_{\ce{N2}} = -1/n_{\ce{N2}} < 0$: siempre en sentido directo.
\end{solution}

\begin{solution}{exo:b2:equilibrium-shifts:11}
$K^\circ = \exp(-1620/(8.314 \times 1100)) = 0.84$: $p(\ce{CO2}) =
\qty{0.84}{bar}$ en el equilibrio, es decir, $n = pV/RT = 0.84 \times 10^5 \times
0.0100/(8.314 \times 1100) = \qty{0.092}{mol}$ de gas. Con \qty{0.050}{mol}
de caliza, se descompone toda ($p = \qty{0.46}{bar} < K^\circ p^\circ$):
no hay equilibrio, hay ruptura de equilibrio. Con \qty{0.200}{mol}, equilibrio:
\qty{0.092}{mol} de \ce{CO2} y de \ce{CaO}, y quedan \qty{0.108}{mol} de \ce{CaCO3}.
\end{solution}

\begin{solution}{exo:b2:equilibrium-shifts:12}
El lecho 1 sale a unos \qty{890}{K} con un 68\,\% de conversión, el lecho 2 a unos
\qty{785}{K} con un 92\,\% y el lecho 3 a unos \qty{725}{K} con un 98\,\%. En un solo
lecho adiabático el gas se calienta al reaccionar hasta alcanzar la curva de
equilibrio, que a esa temperatura solo permite en torno a un 68\,\%. Un catalizador frío
todo el tiempo no es posible porque el catalizador solo funciona en caliente; enfriar
entre lechos mantiene la velocidad y sube el límite de equilibrio. El último lecho
está limitado por el equilibrio a la temperatura más baja a la que el catalizador
aún funciona (las plantas absorben entonces el trióxido y hacen pasar el gas una vez más sobre
un catalizador).
\end{solution}

\begin{solution}{pb:b2:equilibrium-shifts:1}
\textbf{1.} $\Delta_r H^\circ = \qty{-91.88}{kJ/mol}$, $\Delta_r S^\circ =
2(192.77) - 191.61 - 3(130.68) = \qty{-198.1}{J/(K.mol)}$, $\Delta_r G^\circ =
\qty{-32.8}{kJ/mol}$, $K^\circ = \num{5.6e5}$.
\textbf{2.} $K^\circ(700) = \exp(-54\,380/(8.314 \times 700)) = \num{8.8e-5}$;
$K^\circ(800) = \exp(-77\,320/(8.314 \times 800)) = \num{8.9e-6}$.
\textbf{3.} $\Delta_r H^\circ = -R\ln(K_{800}/K_{700})/(1/800 - 1/700)$,
es decir, \qty{-106}{kJ/mol}, más negativa que a \qty{298}{K}.
\textbf{4.} $\ln K^\circ(723) = \ln(8.75 \times 10^{-5}) + 12\,777 \times
(1/723.15 - 1/700)$, con $12\,777 = 106\,230/8.314$: $K^\circ =
\num{4.9e-5}$.
\textbf{5.} $\Delta_r G^\circ(723) = -91.88 + 723.15 \times 0.1981 =
\qty{51.4}{kJ/mol}$, $K^\circ = \num{1.9e-4}$: cuatro veces demasiado grande, porque
$\Delta_r H^\circ$ y $\Delta_r S^\circ$ varían en \qty{425}{K}
($\Delta_r C_p^\circ \approx \qty{-44}{J/(K.mol)}$).
\textbf{6.} \ce{N2} $1 - \xi$, \ce{H2} $3 - 3\xi$, \ce{NH3} $2\xi$, total $4 -
2\xi$.
\textbf{7.} $x = 2\xi/(4 - 2\xi)$, luego $1 - x = (4 - 4\xi)/(4 - 2\xi)$; la
fracción de nitrógeno es $(1 - \xi)/(4 - 2\xi) = (1 - x)/4$ y la de hidrógeno
es tres veces mayor.
\textbf{8.}
\[
  K^\circ = \frac{x^2}{\frac{1-x}{4}\left(\frac{3(1-x)}{4}\right)^3}
  \left(\frac{p^\circ}{p}\right)^2 = \frac{256}{27}\,\frac{x^2}{(1 - x)^4}
  \left(\frac{p^\circ}{p}\right)^2 ;
\]
se toma la raíz cuadrada.
\textbf{9.} $a = (\sqrt{27}/16) \times \sqrt{4.9 \times 10^{-5}} \times 1 =
\num{2.27e-3}$; $x/(1 - x)^2 = a$ da $x = 0.0023$.
\textbf{10.} $a = 0.455$; $ax^2 - (2a + 1)x + a = 0$: $x = 0.25$.
\textbf{11.} $n = 3$, $r = 1$, $s = 1$ (proporción 1 : 3 mantenida por la reacción),
$\varphi = 1$: $v = 2$. Una vez elegidas $T$ y $p$, no queda nada libre.
\textbf{12.} Presión: sentido directo ($\Delta\nu_{\text{gas}} = -2$). Temperatura:
sentido inverso (\omterm{def:b2:reaction-enthalpy:exothermic}{exotérmica}).
\textbf{13.} Un gas inerte a presión total constante rebaja las presiones parciales
de los gases que reaccionan, como lo haría una caída de presión: sentido inverso.
\textbf{14.} $\partial\ln Q/\partial n_{\ce{N2}} = (-1/0.30 + 2)/n_{\text{tot}} < 0$:
$Q$ cae por debajo de $K^\circ$, sentido directo.
\textbf{15.} No: a volumen constante $\partial\ln Q/\partial n_{\ce{N2}} =
-1/n_{\ce{N2}} < 0$ sea cual sea la composición: sentido directo.
\textbf{16.} Casi sin amoníaco, $Q$ está muy por debajo de $K^\circ$ a la entrada:
el gas entra lejos del equilibrio y reacciona en sentido directo.
\textbf{17.} El gas pasa demasiado poco tiempo sobre el catalizador para que se
alcance el equilibrio; un contacto más largo costaría un convertidor mayor y más
caro para una ganancia pequeña.
\textbf{18.} $x = 2\xi/(4 - 2\xi) = 0.18$ da $\xi = \qty{0.305}{mol}$ por
mol de \ce{N2} alimentado: \omterm{def:b2:equilibrium-shifts:single-pass}{conversión por paso} del hidrógeno (y del nitrógeno)
$0.305$, un 31\,\%.
\textbf{19.} En régimen estacionario toda la alimentación fresca acaba convirtiéndose:
salen \qty{2.00}{mol} de amoníaco por unidad de tiempo; por el convertidor deben pasar $1.00/
0.305 = \qty{3.3}{mol}$ de \ce{N2} por unidad de tiempo, y el resto se recicla.
\textbf{20.} El argón (y el metano) que entran con la alimentación nunca reaccionan y
no se retiran con el amoníaco líquido: sin \omterm{def:b2:equilibrium-shifts:single-pass}{purga} se
acumularían en el bucle y rebajarían las presiones parciales de los reactivos.
\textbf{21.} A \qty{500}{K} el catalizador de hierro es demasiado lento: el equilibrio
sería favorable, pero nunca se alcanzaría en un tiempo industrial.
\textbf{22.} A \qty{723}{K} y \qty{200}{bar}, con alimentación estequiométrica y gases
perfectos: $\boldsymbol{x(\ce{NH3}) \approx 0.25}$.
\end{solution}
